\chapter{Tests and the referee subset}\label{ch:tests}

The test suite lives in \texttt{tests/} and runs with \texttt{pytest}. Tests that compare
Exacta with another package skip when the package is missing, and tests that need a GPU skip on
a CPU-only machine; a passing CPU-only run therefore does not establish GPU correctness. The
complete suite takes several minutes on CPU. The \emph{referee subset} selects the tests that
bear on the claims of Chapter~\ref{ch:introduction} and runs in a few minutes on a CPU-only
machine.

\section{An isolated environment}\label{tests:referee}

The commands below build the extension and install the test dependencies into a virtual
environment inside the checkout, so that no other Python environment changes. PyTorch is
installed from its CPU index; use the index for your hardware if you also want the GPU tests.
\begin{lstlisting}[language=bash]
git clone --branch main https://github.com/rodolfo-viana/exacta.git
cd exacta
python -m venv .venv && . .venv/bin/activate
python -m pip install torch --index-url https://download.pytorch.org/whl/cpu
python -m pip install -e ".[test]"
\end{lstlisting}
The editable install compiles the C++ kernels and needs a C++17 compiler (OpenMP is optional).
The reference tools are optional. Without them the tests that use them are skipped, and the
pytest summary counts the skips:
\begin{lstlisting}[language=bash]
python -m pip install monai medpy scikit-image vtk SimpleITK \
  segmentation-models-pytorch \
  "git+https://github.com/google-deepmind/surface-distance.git" \
  "git+https://github.com/gasperpodobnik/MeshMetrics.git"
\end{lstlisting}

\section{Running the referee subset}\label{tests:running}

\begin{lstlisting}[language=bash]
python -m pytest -m referee -q -rs     # the referee subset; -rs lists the skips
python -m pytest -q                    # the complete suite
python examples/driver.py              # the example driver
\end{lstlisting}

\texttt{tests/referee\_subset.txt} lists the tests of the subset as pytest node ids, with
comments that say why each group is there. \texttt{tests/conftest.py} adds the marker
\texttt{referee}, registered in \texttt{pyproject.toml}, to every test that matches an entry:
a file selects all its tests, \texttt{file::test} all parametrizations of a test. No existing
test file is edited. If a listed file is collected but an entry matches nothing, pytest stops
with an error, so a renamed test cannot drop out of the subset unnoticed; the test
\texttt{tests/test\_referee\_subset.py} checks the listing as well.

The subset covers, by group of the listing:
\begin{description}[nosep,leftmargin=0pt,labelindent=0pt]
\sloppy
\item[Exact distance transform.] All of \texttt{test\_edt.py}: agreement with SciPy, equality
  with a brute-force search for the exact floating-point minimum (a property test over random
  masks and spacings), the lines that the exactness argument produces (parabolas that tie far
  apart and within rounding error), extreme spacings, the overflow guard and
  thread-count independence.
\item[Correctly rounded sums.] \texttt{test\_sums.py}: the exact sum equals \texttt{math.fsum}
  in any order, and surface means do not depend on the thread count (on GPU, also on the device).
\item[Metric presets against reference tools.] The surface metrics are compared with NumPy and
  SciPy references and with MONAI, MedPy, DeepMind's \texttt{surface-distance}, MeshMetrics and
  VTK, including NSD ties, in \texttt{test\_surface.py}, \texttt{test\_surfels.py} and
  \texttt{test\_mesh.py}. \texttt{test\_conventions.py} checks the reporting signatures.
  Property tests cover HD, ASSD and NSD, and the overlap and volume scores are compared with MedPy,
  MONAI and SimpleITK.
\item[Losses against compositions.] Independent formulas, MONAI and
  \texttt{segmentation\_models\_pytorch}; fused or composed losses against separately composed
  PyTorch terms, values and gradients, and finite-difference gradient checks.
\item[The example driver.] \texttt{test\_example\_driver.py} (Chapter~\ref{ch:driver}).
\end{description}

\section{Size and runtime}\label{tests:runtime}

On a CPU-only machine with the reference tools of Table~\ref{tab:versions} installed, pytest
collects 2713 tests. The subset selects 614 of them: 581 pass and 33 are skipped because they
need CUDA and Triton, and none is skipped for a missing package. The complete suite runs
2119 tests and skips 594.

Run times, with two OpenMP threads, CUDA hidden, on an AMD Ryzen 9 5900XT (16 cores):
\begin{center}
\small
\begin{tabular}{@{}lr@{}}
\toprule
Run & Test time \\
\midrule
Referee subset, three runs & 27.5, 24.0 and 23.6 s \\
Referee subset, restricted to two cores (\texttt{taskset}) & 25.1 s \\
Referee subset, first measurement under heavy load from other jobs & 70 s \\
Complete suite & 106 s \\
\bottomrule
\end{tabular}
\end{center}
The wall-clock time is about 3 to 5 s longer than the test time reported by pytest (interpreter
start and collection). The numbers depend on the machine, its load and the installed packages;
a slower machine takes proportionally longer. The subset takes about a fifth of the time of the
complete suite because it leaves out the distributed (multi-process), \texttt{torch.compile}
and many input-validation tests, which are slow and do not bear on exactness.

\begin{table}[ht]
\centering
\small
\begin{tabular}{@{}ll@{\qquad}ll@{}}
\toprule
Package & Version & Package & Version \\
\midrule
Python & 3.14.4 & MONAI & 1.6.1 \\
PyTorch & 2.14.1 & MedPy & 0.5.2 \\
NumPy & 2.5.4 & DeepMind \texttt{surface-distance} & 0.1 \\
SciPy & 1.18.1 & MeshMetrics (PyMeshMetrics) & 1.4.1 \\
pytest & 9.1.1 & SimpleITK & 2.5.6 \\
Hypothesis & 6.170.8 & VTK & 9.7.1 \\
scikit-image & 0.26.0 & segmentation-models-pytorch & 0.6.0 \\
\bottomrule
\end{tabular}
\caption{Versions installed when the subset was timed.}
\label{tab:versions}
\end{table}

\section{Checking portability}\label{tests:portability}

\texttt{scripts/portability\_digest.py run} evaluates a fixed set of computations (the
distance transform, connected components, skeletons and the surface metrics with every surface
model) on inputs generated from integer arithmetic, and writes the SHA-256 of each result.
\texttt{compare} checks that the digest files written on different machines agree exactly. This
is the tool behind the statement that the surface scores have the same bits on every platform
and thread count that it was run on.
\begin{lstlisting}[language=bash]
python scripts/portability_digest.py run --out digest-here.json
python scripts/portability_digest.py compare digest-here.json digest-there.json
\end{lstlisting}
